\documentclass[12pt]{article}
\usepackage[utf8]{inputenc}
\usepackage{parskip}
\usepackage[spanish]{babel}
\usepackage{amsmath, amssymb}
\usepackage{graphicx}
\usepackage{makecell}
\usepackage{geometry}
\usepackage{url}
\usepackage[hidelinks]{hyperref}

% --------------------- Comandos para URL y pie de página ---------------------
\makeatletter
\g@addto@macro{\UrlBreaks}{\UrlOrds}
\makeatother

\renewcommand{\footnoterule}{%
  \kern -3pt                         % Espacio vertical sobre la línea
  \hrule width \textwidth height 0.4pt % Ancho total de la página y grosor
  \kern 2.6pt                        % Espacio vertical debajo de la línea
}
% -----------------------------------------------------------------------

\geometry{letterpaper, margin=2.5cm}
\setlength{\parskip}{1em}
\begin{document}
\begin{titlepage}
\begin{center}
    % Logos
    \begin{minipage}{0.3\textwidth}
        \includegraphics[width=0.8\linewidth]{imagenes/UNAM_Logo.png}
    \end{minipage}\hfill
    \begin{minipage}{0.4\textwidth}
        \begin{center}
            \textbf{\Large Universidad Nacional Autónoma de México}\\[0.2cm]
            \textbf{Facultad de Ciencias}\\[0.2cm]
        \textbf{Bases de Datos}
        \end{center}
    \end{minipage}\hfill
    \begin{minipage}{0.3\textwidth}
        \includegraphics[width=0.8\linewidth]{imagenes/Logo_FC.png}
    \end{minipage}
    \vspace{3cm}
    \textbf{\large Practica 3 (Parte 2)}\\

    Licenciatura en Ciencias de la Computación\\
    Profesor: Gerardo Avilés Rosas \\
    \vspace{2cm}
    \textbf{Integrantes: }\\
    Cisneros Tenorio Diego Alexander\\
    Nuñez Ruiz José Manuel \\
    Olmos Ortiz Andrea Lucia\\
    Reynoso Ascencio Alexa \\
    Vega Navas Saúl \\
    \vspace{2cm}
    \textbf{Fecha de entrega: } 01 de octubre del 2026
    \vfill
\end{center}
\end{titlepage}

\tableofcontents
\newpage

\section*{Correcciones respecto a la práctica pasada}
\addcontentsline{toc}{section}{Correcciones respecto a la práctica pasada}

\newpage
\section*{Consideraciones de diseño}
\addcontentsline{toc}{section}{Consideraciones de diseño}

En cuanto al diseño de las entidades:
\begin{itemize}
    \item Se optó por otorgarles identificadores a algunas entidades que consideramos son lo suficientemente importantes como para mantenerse como entidades fuertes: Sucursal, juego, partida, tarjeta, canje, cliente y premio.
    \item Por otro lado, se optó por usar recarga y mantenimiento como entidades débiles pues consideramos que estas entidades están fuertemente ligadas a otras y no pueden existir por sí mismas.
    \item En cuanto a los empleados, se decidió usar herencia, pues en el caso de uso los trabajadores comparten la mayoría de sus atributos y además cada uno tiene al menos uno más por sí mismo. A pesar de que gerente, cajero y encargado de premios tienen exactamente los mismos atributos, cada entidad tiene acciones y formas diferentes de interactuar con el resto de entidades, por lo que consideramos que es mejor seguir el caso de uso y mantenerlas por separado.
    \item En cuanto a los juegos, consideramos mapear tal cual el caso de uso, dividiéndolos en 4 tipos, asignándoles un atributo sobre los puntos que generan cada uno y análogamente a los empleados se decidió mantenerlas por separado pues son distintos juegos a clasificar.
    \item Para todas las entidades que tienen como atributos una fecha, hora o nombre, se decidió hacer un atributo compuesto para garantizar la atomicidad de los atributos.
\end{itemize}

\newpage
\section*{Relaciones fuertes}
\addcontentsline{toc}{section}{Relaciones fuertes}

\begin{enumerate}
    \item \textbf{REGISTRAR} Cliente -- Tarjeta (1:1)
    \begin{itemize}
        \item \textbf{Cliente:} total (1,1). Todo cliente registrado en el sistema debe poseer obligatoriamente una y solo una tarjeta activa para ser identificado.
        \item \textbf{Tarjeta:} total (1,1). Una tarjeta en el sistema debe estar asociada a un único cliente registrado.
        \item Todo cliente requiere de únicamente una tarjeta, por lo que cada tarjeta está asignada a un único cliente.
    \end{itemize}

    \item \textbf{JUGAR} Partida -- Tarjeta (N:1)
    \begin{itemize}
        \item \textbf{Partida:} total (1,M). Una partida requiere ser jugada mediante al menos una tarjeta activa.
        \item \textbf{Tarjeta:} parcial (0,N). Una tarjeta puede no haber jugado una tarjeta todavía o haber jugado múltiples partidas a lo largo del tiempo.
        \item Para iniciar una partida se requiere de una tarjeta activada y se pueden registrar varias tarjetas para una misma partida en multijugador, mientras que una misma tarjeta puede participar en múltiples partidas.
    \end{itemize}

    \item \textbf{REPRODUCIR} Máquina -- Partida (1:N)
    \begin{itemize}
        \item \textbf{Máquina:} parcial (0,N). Una máquina recién instalada o fuera de servicio puede no tener partidas registradas, o haber reproducido muchas partidas si está activa.
        \item \textbf{Partida:} total (1,1). Toda partida debe ejecutarse necesariamente en una máquina física específica.
        \item Una máquina puede generar varias partidas durante un cierto lapso de tiempo, pero cada partida solo se ejecuta en una única máquina.
    \end{itemize}

    \item \textbf{TENER} Sucursal -- Premio (N:M)
    \begin{itemize}
        \item \textbf{Sucursal:} parcial (0,N). Una sucursal puede estar temporalmente sin un premio en particular o tener múltiples premios.
        \item \textbf{Premio:} parcial (0,M). Un tipo de premio puede no estar disponible en una sucursal específica o formar parte del inventario de varias sucursales.
        \item Una sucursal puede almacenar múltiples premios, mientras que un mismo tipo de premio puede estar disponible en múltiples sucursales.
    \end{itemize}

    \item \textbf{PERTENECER} Sucursal -- Máquina (1:N)
    \begin{itemize}
        \item \textbf{Sucursal:} parcial (0,N). Una sucursal en remodelación puede no contar con máquinas asignadas, mientras que una sucursal activa alberga muchas.
        \item \textbf{Máquina:} parcial (0,1). Una máquina física puede estar en traslado, mantenimiento o sin asignación, o estar ubicada en a lo más una sucursal.
        \item Una sucursal puede no tener máquinas, si apenas se está construyendo o remodelando, y una máquina puede pertenecer a lo más en una sucursal o por ejemplo en reparación fuera de una sucursal.
    \end{itemize}

    \item \textbf{INSTALAR} Máquina -- Juego (N:1)
    \begin{itemize}
        \item \textbf{Máquina:} total (1,1). Toda máquina debe de tener instalado obligatoriamente un juego para poder operar.
        \item \textbf{Juego:} parcial (0,N). Un juego puede no estar instalado todavía en ninguna máquina o estar instalado en múltiples máquinas.
        \item Toda máquina debe de tener instalada un juego, mientras que un mismo juego puede estar instalado en varias máquinas o estar en desarrollo.
    \end{itemize}

    \item \textbf{REPARAR} Técnico -- Máquina (N:M)
    \begin{itemize}
        \item \textbf{Técnico:} parcial (0,N). Un técnico contratado puede no haber realizado reparaciones aún o haber reparado múltiples máquinas.
        \item \textbf{Máquina:} parcial (0,M). Una máquina nueva puede no haber necesitado reparaciones o ser reparada a lo largo de su vida útil.
        \item Un técnico puede reparar múltiples máquinas y una máquina puede ser reparada por distintos técnicos en diferentes ocasiones.
    \end{itemize}

    \item \textbf{TRABAJAR} Empleados -- Sucursal (N:M)
    \begin{itemize}
        \item \textbf{Empleados:} total (1,N). Todo empleado debe estar asignado obligatoriamente a al menos una sucursal para trabajar.
        \item \textbf{Sucursal:} total (1,M). Una sucursal debe contar obligatoriamente con un equipo de al menos un empleado para poder operar.
        \item Una sucursal cuenta con múltiples empleados y un empleado debe trabajar en una o más sucursales.
    \end{itemize}

    \item \textbf{ADMINISTRAR} Gerente -- Sucursal (1:1)
    \begin{itemize}
        \item \textbf{Gerente:} parcial (0,1). Un gerente puede no tener una sucursal asignada en un momento dado o administrar a lo más una sucursal.
        \item \textbf{Sucursal:} total (1,1). Toda sucursal requiere tener obligatoriamente un único gerente asignado.
        \item Un gerente puede administrar una única sucursal y una sucursal debe ser administrada por un único gerente.
    \end{itemize}

    \item \textbf{REALIZAR} Cajero -- Recarga (1:N)
    \begin{itemize}
        \item \textbf{Cajero:} parcial (0,1). Un cajero puede no haber realizado recargas aún o realizar una gran cantidad de recargas.
        \item \textbf{Recarga:} total (1,N). Toda recarga tiene que ser procesada y cobrada obligatoriamente por un único cajero.
        \item Un cajero puede hacer muchas recargas, mientras que una recarga debe ser hecha por un único cajero.
    \end{itemize}

    \item \textbf{EFECTUAR} Tarjeta -- Canje (1:N)
    \begin{itemize}
        \item \textbf{Tarjeta:} parcial (0,N). Una tarjeta puede no haber realizado ningún canje de premios aún o haber acumulado múltiples canjes.
        \item \textbf{Canje:} total (1,1). Todo canje registrado debe estar asociado obligatoriamente a la tarjeta que utilizó sus puntos.
        \item Una tarjeta puede realizar múltiples canjes a lo largo del tiempo, mientras que un transacción de canje pertenece a una sola tarjeta.
    \end{itemize}

    \item \textbf{ATENDER} EncargadoPremios -- Canje (1:N)
    \begin{itemize}
        \item \textbf{EncargadoPremios:} parcial (0,N). Un encargado de premios puede no haber procesado ningún canje aún o haber atendido múltiples solicitudes.
        \item \textbf{Canje:} total (1,1). Todo canje de premios debe ser procesado y entregado obligatoriamente por un único encargado de premios.
        \item Un encargado de premios puede procesar varios canjes durante su jornada, pero cada canje es atendido por un solo encargado.
    \end{itemize}
    
    \item \textbf{OCURRIR\_EN} Sucursal -- Canje (1:N)
    \begin{itemize}
        \item \textbf{Sucursal:} parcial (0,N). Una sucursal nueva puede no registrar canjes aún o llevar a cabo una gran cantidad de ellos.
        \item \textbf{Canje:} total (1,1). Todo canje debe realizarse obligatoriamente dentro de las instalaciones de una sucursal específica.
        \item En una sucursal se pueden llevar a cabo múltiples canjes de premios, mientras que cada transacción de canje ocurre en una única sucursal.
    \end{itemize}

    \item \textbf{INCLUIR} Canje -- Premio (M:N)
    \begin{itemize}
        \item \textbf{Canje:} total (1,N). Todo canje registrado debe incluir obligatoriamente al menos un premio entregado al cliente.
        \item \textbf{Premio:} parcial (0,M). Un premio del catálogo puede no haber sido reclamado todavía en ningún canje o haber sido incluido en múltiples canjes.
        \item Un canje puede involucrar la entrega de uno o varios premios, y un mismo tipo de premio puede ser parte de múltiples canjes distintos.
    \end{itemize}

    \item \textbf{REGISTRAR\_EN} Sucursal -- Recarga (1:N)
    \begin{itemize}
        \item \textbf{Sucursal:} parcial (0,N). Una sucursal puede estar recién abierta y no registrar recargas aún, o haber procesado múltiples recargas.
        \item \textbf{Recarga:} total (1,1). Toda recarga de saldo debe registrarse obligatoriamente en la sucursal física donde fue efectuada.
        \item Una sucursal puede procesar numerosas recargas de tarjetas, mientras que cada operación de recarga se efectúa en una sola sucursal.
    \end{itemize}
    
\end{enumerate}

\newpage
\section*{Relaciones débiles}
\addcontentsline{toc}{section}{Relaciones débiles}

\begin{enumerate}
    \item \textbf{PEDIR} Tarjeta -- Recarga (1:N)
    \begin{itemize}
        \item \textbf{Tarjeta:} parcial (0,N). Una tarjeta recién adquirida puede no haber realizado recargas o contar con múltiples recargas registradas.
        \item \textbf{Recarga:} total (1,1). Una recarga no puede existir sin estar ligada a la tarjeta que recibe el saldo.
        \item Una tarjeta puede recibir múltiples recargas, mientras que una recarga pertenece a una única tarjeta.
    \end{itemize}
\end{enumerate}



\end{document}
